\chapter{Duality and the Global Structure of Conformal Data}
\label{ch:volume-iii-duality-global-structure-conformal-data}

\section{Isomorphism of Theories}

Let \(\mathfrak D_1\) and \(\mathfrak D_2\) be complete conformal data in the
same spacetime dimension.  A duality between them is an isomorphism of
operator algebras and Hilbert spaces that preserves all correlation functions,
Ward identities, and positivity structures after the operators and background
fields have been relabeled.

More explicitly, let \(\Hilb_1(S^{D-1})\) and \(\Hilb_2(S^{D-1})\) be the
radial-quantization Hilbert spaces.  A duality contains a unitary map
\[
  U:\Hilb_1(S^{D-1})\longrightarrow \Hilb_2(S^{D-1})
\]
intertwining the conformal generators:
\[
  U\,Q_A^{(1)}\,U^{-1}=Q_A^{(2)}
\]
for every generator \(Q_A\) of the conformal algebra.  If
\(\mathcal O_i^{(1)}\) is a primary in the first description, then
\[
  U\,\mathcal O_i^{(1)}(x)\,U^{-1}
  =
  M_i{}^j\,\mathcal O_j^{(2)}(x)
\]
inside the space of primaries with the same transformed quantum numbers.  The
matrix \(M_i{}^j\) may mix degenerate primaries.

The equality of local CFT data is then
\begin{equation}
  \Delta_i^{(1)}=\Delta_{M(i)}^{(2)},
  \qquad
  C_{ijk}^{(1)}
  =
  M_i{}^{i'}M_j{}^{j'}M_k{}^{k'}C_{i'j'k'}^{(2)}
  \label{eq:duality-cft-data-isomorphism}
\end{equation}
with the evident extension to spinning structures, global-symmetry
representations, and higher-point functions.  Defect operators, line
operators, and background fields are part of the same complete data when they
are included in the theory.

\begin{definition}[Duality]
  A duality is an isomorphism between complete quantum field theoretic data.
  For a conformal theory, this includes the Hilbert space, conformal symmetry
  action, local operator algebra, allowed extended operators, symmetry
  backgrounds, and all correlation functions, modulo relabeling of equivalent
  bases.
\end{definition}

This definition makes duality a structural statement.  It can be discovered
through two different Lagrangian descriptions, through a geometric
construction, through a lattice regularization, or through a bootstrap
solution.  Once established, it identifies the resulting conformal data.

\section{Quotients of Conformal Manifolds}

Let \(\widetilde{\mathcal M}_{\rm conf}\) be a coordinate cover of a
conformal manifold, and let \(\Gamma\) be a group acting on it by
transformations
\[
  \gamma:\widetilde{\mathcal M}_{\rm conf}
  \longrightarrow
  \widetilde{\mathcal M}_{\rm conf},
  \qquad
  \gamma\in\Gamma,
\]
such that \(\mathfrak D(p)\) and \(\mathfrak D(\gamma p)\) are dual
conformal data.  Then the physical conformal manifold is the quotient
\[
  \mathcal M_{\rm conf}
  =
  \widetilde{\mathcal M}_{\rm conf}/\Gamma,
\]
possibly with orbifold points at fixed loci where additional automorphisms of
the theory appear.

\begin{figure}[H]
  \centering
  \resizebox{0.94\textwidth}{!}{%
  \begin{tikzpicture}[x=1cm,y=1cm,every node/.style={font=\scriptsize}]
    \tikzset{
      arrow/.style={-{Latex[length=2mm]},line width=0.85pt},
      curve/.style={blue!65!black,line width=1.0pt},
      identify/.style={orange!85!black,line width=0.9pt,dashed}
    }
    \begin{scope}[shift={(-3.9,0)}]
      \draw[->,thick] (-1.2,-1.05)--(1.35,-1.05) node[right] {\(\operatorname{Re}\tau\)};
      \draw[->,thick] (-1.2,-1.05)--(-1.2,1.25) node[above] {\(\operatorname{Im}\tau\)};
      \draw[curve] (-0.85,-0.70) .. controls (-0.2,0.40) and (0.42,0.42) .. (0.98,-0.70);
      \draw[identify] (-0.85,-0.70)--(-0.85,1.10);
      \draw[identify] (0.98,-0.70)--(0.98,1.10);
      \node[blue!65!black] at (0.08,1.38) {coordinate cover};
      \node[orange!85!black] at (0.08,-1.55) {\(\tau\sim\gamma\tau\)};
    \end{scope}
    \draw[arrow] (-1.55,0)--(-0.45,0);
    \node[align=center] at (-1.00,0.52) {identify\\dual data};
    \begin{scope}[shift={(2.7,0)}]
      \draw[curve,fill=blue!5] (-1.35,-0.75)
        .. controls (-0.65,0.85) and (0.65,0.85) .. (1.35,-0.75)
        -- cycle;
      \node[violet!80!black] at (-0.62,0.04) {\(\mathfrak D(p)\)};
      \node[violet!80!black] at (0.63,0.04) {\(\mathfrak D(\gamma p)\)};
      \draw[identify] (-0.55,-0.12)--(0.55,-0.12);
      \node[blue!65!black] at (0,-1.25) {quotient conformal manifold};
    \end{scope}
  \end{tikzpicture}
  }
  \caption{A duality group identifies coordinate values whose complete
  conformal data are isomorphic.  The quotient is the global conformal
  manifold.}
  \label{fig:duality-quotient-conformal-manifold}
\end{figure}

For four-dimensional \(\mathcal N=4\) super-Yang--Mills with gauge algebra
\(\mathfrak{su}(N)\), the local coordinate is
\[
  \tau_{\rm YM}
  =
  \frac{\theta}{2\pi}
  +
  \frac{4\pi\ii}{g_{\rm YM}^2},
  \qquad
  \Im\tau_{\rm YM}>0.
\]
The transformations
\begin{equation}
  \tau_{\rm YM}
  \longmapsto
  \frac{a\tau_{\rm YM}+b}{c\tau_{\rm YM}+d},
  \qquad
  a,b,c,d\in\mathbb Z,
  \qquad
  ad-bc=1,
  \label{eq:sl2z-tau-action}
\end{equation}
act by electric-magnetic duality in the standard simply laced cases.  The
precise duality group also depends on the global form of the gauge group and
on the spectrum of line operators.  Thus a statement about
\(\tau_{\rm YM}\) alone is local; the global statement includes line
operators and background fields.

\section{Local Operators and Extended Operators}

Local correlation functions determine much of a CFT, but in a gauge theory
the full quantum field theoretic data also includes extended probes.  A Wilson
line in representation \(R\) of the gauge group is
\[
  W_R(C)
  =
  \operatorname{Tr}_R
  P\exp\left(\ii\int_C A\right),
\]
where \(C\) is an oriented curve and \(P\) denotes path ordering.  A
't Hooft line is specified by a magnetic singularity, equivalently by a
cocharacter of the gauge group modulo Weyl transformations.  A Wilson-'t
Hooft line has both electric and magnetic charges.

Under the transformation \eqref{eq:sl2z-tau-action}, electric and magnetic
charges transform as a rank-two lattice.  If a line operator has charge column
vector
\[
  \begin{pmatrix} q_{\rm e}\\ q_{\rm m}\end{pmatrix},
\]
then the duality transformation acts by an integral linear map compatible with
the Dirac pairing.  The full duality identification therefore preserves the
lattice of allowed line operators, their operator products, and their
correlation functions with local insertions.

The same principle applies to global symmetries.  If \(A_{\rm bg}\) is a
background gauge field for a global symmetry, the generating functional
\[
  Z[g_{\mu\nu},A_{\rm bg},\lambda]
\]
is part of the data.  A duality may act nontrivially on \(A_{\rm bg}\), on
counterterms for background fields, and on the discrete theta angles that
specify the theory.  Anomalies are preserved by the duality map.

\section{Singular Loci and Extra Massless Data}

A conformal manifold can have special loci at which the automorphism group of
the theory is enhanced or at which a chosen Lagrangian coordinate becomes
singular.  At such a point, the quotient
\(\widetilde{\mathcal M}_{\rm conf}/\Gamma\) may have an orbifold
singularity.  The local CFT remains described by conformal data; the
singularity is a statement about the coordinates used to parametrize the
family and about extra identifications of operator bundles.

Operator dimensions may cross along subloci.  When this happens, the correct
object is the vector bundle of degenerate operator spaces, together with the
connection and OPE tensors.  A labeling by individual operators can be chosen
locally after diagonalizing the dilatation operator, but the global
monodromy around a singular locus may return the basis after a nontrivial
linear transformation.

\begin{figure}[H]
  \centering
  \resizebox{0.90\textwidth}{!}{%
  \begin{tikzpicture}[x=1cm,y=1cm,every node/.style={font=\scriptsize}]
    \tikzset{
      arrow/.style={-{Latex[length=2mm]},line width=0.85pt},
      bundle/.style={green!45!black,line width=0.9pt},
      path/.style={orange!85!black,line width=1.0pt,-{Latex[length=2mm]}},
      box/.style={draw,rounded corners=4pt,fill=blue!4,inner sep=4pt,align=center}
    }
    \draw[blue!65!black,line width=1.0pt] (-3.3,0)
      .. controls (-2.0,0.95) and (-0.6,-0.95) .. (0.7,0)
      .. controls (1.75,0.72) and (2.65,0.55) .. (3.3,-0.18);
    \fill[violet!80!black] (-0.02,-0.10) circle (2.4pt);
    \node[violet!80!black] at (-0.02,-0.50) {special locus};
    \draw[path] (-1.25,0.45) arc (150:-170:0.92 and 0.52);
    \foreach \x/\y in {-2.55/0.48,-1.65/0.22,1.55/0.35,2.45/0.08} {
      \draw[bundle] (\x,\y)--++(0,0.72);
      \draw[bundle] (\x,\y)--++(0,-0.72);
    }
    \node[box] at (4.25,0.65) {operator\\bundle};
    \draw[arrow] (3.35,0.20)--(3.80,0.50);
    \node[orange!85!black,align=center] at (-1.55,1.35)
      {transport around the locus\\can mix degenerate primaries};
  \end{tikzpicture}
  }
  \caption{Operator labels live in bundles over the conformal manifold.
  Around special loci, parallel transport can act by a nontrivial monodromy on
  degenerate operator spaces.}
  \label{fig:operator-bundle-monodromy}
\end{figure}

\section{Duality as a Constraint}

Because duality identifies complete conformal data, it turns a local
calculation in one region of \(\mathcal M_{\rm conf}\) into a constraint on
the entire quotient.  For a scalar primary whose transformed label is
\(\gamma\cdot i\),
\[
  \Delta_i(p)
  =
  \Delta_{\gamma\cdot i}(\gamma p),
\]
and similarly for OPE coefficients.  If an operator is invariant under a
duality element up to a phase, its correlation functions transform as modular
or automorphic functions with weights fixed by that phase and by the
normalization of the operator.

The constraint is strongest when combined with analyticity, positivity, and
OPE convergence.  Perturbation theory near one cusp, semiclassical expansion
near another region, and protected supersymmetric data can then determine
nontrivial functions over the conformal manifold.  The result is a global
description of a family of conformal theories rather than a collection of
separate coordinate expansions.
